
\section{Related Work}

\subsection{Learning Rate Schedule}
Learning rate (LR) scheduling is critical for training performant models~\citep{Rethinking,Akhilesh}. Classic schedules like Cosine~\citep{kaplan2020scaling,Simple} or Linear~\citep{defazio2023optimal} decay adjust the LR based on a predefined total training duration, which is inflexible for continual training. The Warmup-Stable-Decay (WSD) schedule~\citep{minicpm} addresses this by introducing a stable LR phase after warmup, decoupling the eventual decay from a fixed training length and offering greater flexibility for long or continuous training runs. More recently, researches have explored ``schedule-free'' methods to eliminate the decay phase entirely, which maintain a constant LR and instead leverage weight averaging techniques~\citep{Aaron2024Road,song2025through,Hanlin2025cbs}. Builds upon these schedule-free principles, our work propose to replace WSD's decay phase with a post-hoc checkpoint merging operation instead of specific online averaging strategies. This simplifies the training pipeline and, by formalizing the connection between LR decay and checkpoint merging, allows decay strategy to be theoretically approximated, leading to free offline exploration and enhanced model performance.

\subsection{Model Merging}
Model merging~\citep{Averaging,Modelsoups} has emerged as an efficient paradigm for model construction. This approach achieves effective knowledge transfer and performance improvement through parameter-level integration of multiple models.
Model merging is primarily utilized in two distinct scenarios: (1) The integration of knowledge and capabilities from multiple independently trained models into a single parameter set, with the objective of preserving maximal performance from each specialized model~\citep{commanda,WARP};  and (2) the merging of checkpoints along a single training trajectory, which functions as a form of Polyak averaging~\citep{polyak1992acceleration}. In this application, the merging operation acts as a smoothing mechanism that reduces noise inherent in stochastic gradient-based optimization~\citep{sanyal2023early,Checkpoint,lawa,Trainable,Sandler2023training,1sqrt}. While \cite{pma} show model merging can achieve
performance competitive with decay-based schedule, these techniques have also demonstrated practical utility in industrial-scale LLM development~\citep{grattafiori2024llama,deepseekv3,commanda}.
Our work falls into the second category of merging checkpoints along a single trajectory. We establish a theoretical connection between this operation and learning rate decay and provide a principled approach to convert various
LR decay strategies into a theoretically approximate model averaging implementation. 
